\section{The Proportionality Compromise}
\label{sec:compromise}

The central empirical question is: what is the minimum partisan information
required to close the proportionality gap in each state?

\subsection{The Compromise Spectrum}

We define three categories:

\paragraph{Free proportionality states.}
States where $\sigma < 0.05$ (partisan signal $< 5\,\%$ shift from neutral
distribution). In these states, geography-only algorithms already produce
approximately proportional outcomes. No partisan constraint is needed.
These are states with relatively uniform partisan geography.

\paragraph{Cheap proportionality states.}
States where $0.05 \leq \sigma < 0.15$. A modest partisan constraint
(at $\eta \approx 1.10$) closes the proportionality gap.
These states have moderate partisan concentration but sufficient geographic
mixing that the R-bloc can be constructed from majority-R tracts without
requiring extreme packing.

\paragraph{Expensive proportionality states.}
States where $\sigma \geq 0.15$ or where Lorenz feasibility is violated.
Closing the proportionality gap requires either a strong partisan constraint
(high information input) or is geometrically impossible due to extreme
concentration. Georgia and Massachusetts are the clearest examples:
Atlanta and Boston are so dominant that any contiguous R-bloc contains
a majority of D voters, making proportional R districts unachievable
without non-contiguous or extremely gerrymandered maps.

\subsection{The Inverse Relationship with the Rodden Effect}

A striking pattern emerges: the states where proportionality is most expensive
are exactly the states with the largest Rodden gaps (Table~\ref{tab:gap-baseline}).
This is not a coincidence. Both phenomena have the same root cause: extreme
Democratic geographic concentration.

\begin{proposition}
\label{prop:rodden-sigma}
The partisan signal strength $\sigma$ and the Rodden proportionality gap are
both increasing functions of the partisan Lorenz curve's deviation from the
diagonal. States with the largest B.9 proportionality gaps require the most
partisan information to achieve proportionality.
\end{proposition}

This has a direct implication for redistricting reform: the states that most
need proportional outcomes (where geographic sorting creates the largest gaps)
are exactly the states where achieving proportionality requires the most
politically controversial partisan inputs. There is no easy geographic fix
for extreme partisan concentration.

\subsection{The Wisconsin Case}

Wisconsin ($d = 0.503$, $k=8$) is a borderline case in two senses.
First, it is exactly at the proportional threshold: $k_D = 4$ is proportional
for $d = 0.503$.
Second, the Lorenz analysis places it in the ``moderate partisan concentration''
category: $\sigma \approx 0.08$.

Under B.12 at $\eta = 1.10$, Wisconsin should achieve approximately 4D/4R ---
proportional for a near-even state.
Under B.9 (geography only), it produces 3D/5R ($-12.8$\,pp gap).
The difference of 1 seat over 8 corresponds to a $12.5$ percentage-point
change in Democratic representation, achieved by a partisan constraint
of $\sigma = 0.08$ (shifting 8\,\% of D votes from the neutral allocation).

This illustrates the core tradeoff: a modest partisan signal (8\,\% reallocation)
produces a large proportionality improvement (12.5\,pp), but requires
explicitly incorporating partisan vote data into the redistricting algorithm.

\subsection{Sensitivity to Geographic Resolution}
\label{sec:compromise:maup}

We repeated the 6-state empirical table at block-group resolution
($n \approx 239{,}000$ units nationally, compared to $\approx 74{,}000$ census
tracts).
Results are qualitatively identical: Wisconsin improves at $\eta = 1.10$
(achieving near-proportional outcomes as at tract resolution), and North Carolina
worsens at all $\eta$ values (the B.12 constraint increases the proportionality
gap regardless of resolution).
The proportionality paradox --- competitive states showing large Rodden gaps
despite near-zero $\sigma$ --- is not a resolution artifact; it reflects the
geographic distribution of partisan populations, which is stable across
geographic scales.
This finding is consistent with prior MAUP literature
\citep{chen2013} showing that partisan sorting patterns are robust to
unit aggregation.

\subsection{The Recursive Case}

Theorem~\ref{thm:proportionality} guarantees proportionality at the first
bisection level.
For full proportionality across the entire $k$-district map, the same
constraint must be applied recursively at every level of the PrimeFactor
bisection tree.

At each level, the local Democratic fraction $d'$ of the subregion determines
a new $\mathtt{tpwgts}$ vector.
If the D-bloc is a majority-D subregion ($d' > 0.5$), the recursive constraint
uses the local $d'$ and $k' = k_D$ to produce $k'_D$ and $k'_R$ for that level.

The full recursive proportionality guarantee follows by induction:
each bisection achieves local proportionality, and the leaf districts
each contain one winning party at the correct population level.
